\introduction
\markboth{Introduction générale}{Introduction générale}

Le contrôle d'accès vise à réserver l'entrée d'un lieu aux personnes autorisées. Il repose classiquement sur un objet détenu, tel qu'une clé ou un badge, ou sur un secret, tel qu'un code. Ces facteurs authentifient la possession d'un objet ou la connaissance d'une information, et non la personne elle-même : ils peuvent être copiés, prêtés, perdus ou observés. Les systèmes à badge ou à code conservent un historique des passages, mais celui-ci identifie le badge ou le code utilisé, et non son porteur.

La biométrie remplace ces facteurs par une caractéristique propre à l'individu. La reconnaissance faciale a pour avantage de s'effectuer sans contact, à partir d'une simple caméra. Elle repose cependant sur une comparaison statistique dont le résultat est une mesure de ressemblance et non une certitude. La décision d'accès est donc prise sous incertitude et expose à deux types d'erreur : la fausse acceptation, qui admet une personne non autorisée, et le faux rejet, qui refuse une personne autorisée~\cite{nist_far,nist_frr}.

\section*{1- Contexte et justification}

Une architecture courante de reconnaissance faciale consiste à transmettre les images à des serveurs distants, où elles sont analysées. Ce modèle est désigné par le terme d'informatique en nuage, ou \gls{cloudcomputing}. Les caméras connectées de la marque Ring en sont un exemple : en 2023, la \gls{ftc} américaine a relevé que leurs vidéos étaient stockées sans chiffrement sur le réseau de l'entreprise et accessibles à tous ses employés~\cite{ftc2023ring}. Ce modèle suppose une alimentation électrique et un accès à Internet continus. Au Bénin, le coût d'Internet et l'instabilité de l'alimentation électrique rendent ces conditions difficiles à garantir, d'après l'observation de l'auteur. Lorsqu'elles ne le sont pas, le service s'interrompt avec la liaison.

Les données biométriques font par ailleurs l'objet d'un encadrement juridique particulier. Le \gls{rgpd} range les données biométriques utilisées pour identifier une personne de façon unique parmi les catégories particulières de données, dont le traitement est en principe interdit sauf exceptions~\cite{ue2016rgpd}. Au Bénin, le Code du numérique pose la même interdiction de principe à son article 394~\cite{apdp2018codenumerique}. Contrairement à un mot de passe, un visage ne peut pas être remplacé après une fuite, ce qui justifie de limiter sa transmission et sa conservation.

Le traitement en périphérie de réseau, ou \gls{edge}~\cite{shi2016edge}, consiste à analyser les données à proximité de leur source plutôt que dans un centre de données distant. Il réduit la dépendance à la connexion et permet de conserver les images sur le site. Il impose en contrepartie de travailler avec des ressources de calcul limitées.

\section*{2- Problématique}

Le contrôle d'accès par reconnaissance faciale pose plusieurs problèmes.

Sur le plan technique, l'analyse des images demande beaucoup de calculs, alors qu'un équipement installé sur place dispose de peu de ressources.

Sur le plan de la sécurité, le système peut être trompé par une photographie, une vidéo ou un écran présenté devant la caméra~\cite{chakraborty2014overview,iso30107}.

Enfin, le visage d'un utilisateur change avec le temps, ce qui dégrade la reconnaissance si les références enregistrées ne sont pas mises à jour~\cite{bestrowden2018longitudinal,pisani2019adaptive}. Or une mise à jour sans contrôle peut introduire des erreurs dans les références, comme l'a montré une première version de BioGate.

La problématique centrale de ce mémoire peut donc être formulée comme suit : \textit{comment concevoir un système de contrôle d'accès par reconnaissance faciale, exécuté en local selon le principe de l'Edge Computing, qui résiste aux attaques par présentation et reste fiable malgré l'évolution du visage de ses utilisateurs ?}

\section*{3- Objectifs de la recherche}

L'objectif général est de concevoir et de réaliser \textbf{BioGate}, un système de contrôle d'accès par reconnaissance faciale fonctionnant en périphérie de réseau.

Les objectifs spécifiques sont les suivants :
\begin{itemize}
  \item concevoir une chaîne de traitement d'images adaptée aux ressources d'un équipement local ;
  \item définir une règle de décision fondée sur un score de confiance gradué (\gls{trustscore}), qui associe la reconnaissance du visage à un second facteur d'authentification ;
  \item intégrer une détection de vivacité (\gls{liveness}) pour limiter les attaques par présentation ;
  \item encadrer la mise à jour des références biométriques par une confirmation indépendante du visage ;
  \item limiter la transmission et la conservation des données biométriques ;
  \item évaluer le système par des mesures : taux de fausses acceptations et de faux rejets, résistance aux attaques par présentation, temps de traitement et comportement en cas de panne.
\end{itemize}

\section*{4- Organisation du mémoire}

Ce document s'organise en trois chapitres :
\begin{itemize}
  \item le \textbf{chapitre~\ref{chap:1}} présente l'état de l'art : les systèmes de contrôle d'accès et leurs limites, le traitement en périphérie, les bases de la reconnaissance de visages et de la détection de vivacité, ainsi que le cadre éthique et légal ;
  \item le \textbf{chapitre~\ref{chap:2}} analyse les besoins et décrit la conception du système : architecture, chaîne de traitement, règle de décision et mécanismes contre l'usurpation ;
  \item le \textbf{chapitre~\ref{chap:3}} décrit la réalisation, puis l'évaluation chiffrée du système, ses limites et ses perspectives.
\end{itemize}
